%------------------------------------------------------------------------------%
\begin{question}
  Calcule o vetor gradiente da função \(f(x, y) = \sqrt{2x+3y}\)
  no ponto \((-1, 2)\)
\end{question}

%------------------------------------------------------------------------------%
\begin{resolution}{Solução}
  Calculando as derivadas parciais

  \pause
  \[
    \D{f}{x}
    \pause
    \;=\; \D{}{x}\sqrt{2x+3y}
    \pause
    \;=\; \D{}{x}\left(2x+3y\right)^{\sfrac{1}{2}}
  \]
  \[
    \hphantom{\D{f}{x}}
    \pause
    \;=\; \frac{1}{2}\left(2x+3y\right)^{\sfrac{-1}{2}}\D{}{x}\left(2x+3y\right)
    \pause
    \;=\; \frac{1}{\sqrt{2x+3y}}
  \]

  \pause
  \[
    \D{f}{x}(-1, 2)
    \pause
    \;=\; \frac{1}{\sqrt{2(-1)+3(2)}}
    \pause
    \;=\; \frac{1}{\sqrt{-2+6}}
    \pause
    \;=\; \frac{1}{\sqrt{4}}
    \pause
    \;=\; \frac{1}{2}
  \]
\end{resolution}

%------------------------------------------------------------------------------%
\begin{resolution}{Solução}
  Calculando as derivadas parciais

  \pause
  \[
    \D{f}{y}
    \pause
    \;=\; \D{}{y}\sqrt{2x+3y}
    \pause
    \;=\; \D{}{y}\left(2x+3y\right)^{\sfrac{1}{2}}
  \]
  \[
    \hphantom{\D{f}{y}}
    \pause
    \;=\; \frac{1}{2}\left(2x+3y\right)^{\sfrac{-1}{2}}\D{}{y}\left(2x+3y\right)
    \pause
    \;=\; \frac{3}{2\sqrt{2x+3y}}
  \]

  \pause
  \[
    \D{f}{y}(-1, 2)
    \pause
    \;=\; \frac{3}{2\sqrt{2(-1)+3(2)}}
    \pause
    \;=\; \frac{3}{2\sqrt{-2+6}}
    \pause
    \;=\; \frac{3}{2\sqrt{4}}
    \pause
    \;=\; \frac{3}{4}
  \]
\end{resolution}

%------------------------------------------------------------------------------%
\begin{resolution}{Solução}
  O vetor gradiente
  \[
    \nabla f
    \pause
    \;=\; \left(\begin{array}{c}
        \displaystyle\D{f}{x} \\[5mm]
        \displaystyle\D{f}{y}
      \end{array}\right)
    \pause
    \;=\; \left(\begin{array}{c}
             \dfrac{1}{\sqrt{2x+3y}} \\[5mm]
             \dfrac{3}{2\sqrt{2x+3y}}
          \end{array}\right)
    \pause
    \;=\; \frac{1}{\sqrt{2x+3y}} \left(\begin{array}{c}
            1 \\[2mm]
            \sfrac{3}{2}
          \end{array}\right)
  \]
\end{resolution}

%------------------------------------------------------------------------------%
\begin{resolution}{Solução}
  O vetor gradiente no ponto \((-1, 2)\)
  \[
    \nabla f(-1, 2)
    \;=\; \left[\frac{1}{\sqrt{2x+3y}} \left(\begin{array}{c}
      1 \\[2mm]
      \sfrac{3}{2}
    \end{array}\right)\right]\evalat{(-1, 2)}{}
    \;=\; \vetor{\sfrac{1}{2}}{\sfrac{3}{4}}
  \]
\end{resolution}

%------------------------------------------------------------------------------%
